% !TeX program = xelatex
% ============================================================================
% 论文中文译文模板（单栏 / Elsevier 风格）
% 用法：
%   1. 用 scripts/probe_layout.py 探测原 PDF，把下面的版式参数填进来
%   2. 替换标题区与正文
%   3. 用 scripts/build_check.py 编译校验
% ============================================================================
\documentclass[10pt]{ctexart}

% ---------- 版式参数：按 probe_layout.py 的输出填写 ----------
\newcommand{\PageW}{544.25pt}      % 页面宽
\newcommand{\PageH}{742.68pt}      % 页面高
\newcommand{\MargL}{39pt}          % 左边距
\newcommand{\MargR}{42.5pt}        % 右边距
\newcommand{\MargT}{64pt}          % 上边距（含页眉区）
\newcommand{\MargB}{38pt}          % 下边距
\newcommand{\HeadL}{Author et al.} % 页眉左
\newcommand{\HeadR}{Journal Name (2026) 000000} % 页眉右

\usepackage[
  paperwidth=\PageW, paperheight=\PageH,
  left=\MargL, right=\MargR, top=\MargT, bottom=\MargB,
  headheight=12pt, headsep=16pt, footskip=26pt
]{geometry}

% ---------- 字体：中文宋/黑/楷 + 西文 Times + 公式 STIX ----------
\usepackage{fontspec}
\setmainfont{Times New Roman}
\setsansfont{Arial}
\setCJKmainfont{SimSun}
\setCJKsansfont{SimHei}
\setCJKfamilyfont{kai}{KaiTi}
\newcommand{\kai}{\CJKfamily{kai}}
\usepackage{unicode-math}
\setmathfont{STIX Two Math}

% ---------- 常用宏包 ----------
\usepackage{graphicx}
\usepackage{amsmath}
\usepackage{booktabs}
\usepackage{multirow}
\usepackage{array}
\usepackage{caption}
\usepackage{float}
\usepackage{enumitem}
\usepackage{xcolor}
\usepackage{fancyhdr}
\usepackage{algorithm}
\usepackage{algpseudocode}
\usepackage[hidelinks]{hyperref}

\captionsetup{font=small, labelsep=quad}
\setlist{nosep, leftmargin=1.6em}
\floatname{algorithm}{算法}
\renewcommand{\algorithmicrequire}{\textbf{输入：}}
\renewcommand{\algorithmicensure}{\textbf{输出：}}

% 行距：原文字号 9pt/行距 10.5pt 时，中文用 1.30 视觉密度接近
\linespread{1.30}

% ---------- 页眉页脚 ----------
\pagestyle{fancy}
\fancyhf{}
\fancyhead[L]{\footnotesize\itshape \HeadL}
\fancyhead[R]{\footnotesize\itshape \HeadR}
\fancyfoot[C]{\footnotesize\thepage}
\renewcommand{\headrulewidth}{0pt}
\fancypagestyle{plain}{%
  \fancyhf{}%
  \fancyhead[L]{\footnotesize\itshape \HeadL}%
  \fancyhead[R]{\footnotesize\itshape \HeadR}%
  \fancyfoot[C]{\footnotesize\thepage}%
  \renewcommand{\headrulewidth}{0pt}%
}

% ---------- 标题层级（仿期刊：黑体、左对齐、编号 1 / 1.1 / 1.1.1） ----------
\ctexset{
  section/format       = \raggedright\heiti\fontsize{11.5pt}{14pt}\selectfont,
  section/aftername    = \quad,
  section/beforeskip   = 14pt plus 2pt minus 2pt,
  section/afterskip    = 8pt plus 1pt,
  subsection/format    = \raggedright\heiti\fontsize{10.5pt}{13pt}\selectfont,
  subsection/aftername = \quad,
  subsection/beforeskip= 10pt plus 1pt,
  subsection/afterskip = 6pt,
  subsubsection/format = \raggedright\heiti\fontsize{10pt}{12.5pt}\selectfont,
  subsubsection/aftername = \quad,
  subsubsection/beforeskip = 8pt,
  subsubsection/afterskip  = 5pt,
}
\setlength{\parindent}{1.2em}
\setlength{\parskip}{0pt}

\begin{document}

% ================= 标题区 =================
\begin{flushleft}
  {\heiti\fontsize{20pt}{25pt}\selectfont
   中文标题第一行：副标题前的主标题\\[2pt]
   中文标题第二行\par}
\end{flushleft}

\vspace{8pt}
{\noindent\fontsize{11pt}{14pt}\selectfont
作者一\textsuperscript{a}，作者二\textsuperscript{a}，作者三\textsuperscript{b,*}\par}

\vspace{6pt}
{\noindent\footnotesize\kai\raggedright
\textsuperscript{a}~单位一，城市；\quad
\textsuperscript{b}~单位二，城市\par}

\vspace{4pt}
\hrule height 0.4pt
\vspace{6pt}

% ================= 关键词 + 摘要（双栏） =================
\noindent
\begin{minipage}[t]{0.26\textwidth}
  \vspace{0pt}
  {\footnotesize\heiti 文章信息\par}
  \vspace{3pt}\hrule height 0.3pt\vspace{4pt}
  {\footnotesize\heiti 关键词：\par}
  \vspace{2pt}
  {\footnotesize 关键词一\\关键词二\\关键词三\par}
\end{minipage}%
\hfill
\begin{minipage}[t]{0.70\textwidth}
  \vspace{0pt}
  {\footnotesize\heiti 摘\quad 要\par}
  \vspace{3pt}\hrule height 0.3pt\vspace{4pt}
  {\footnotesize 摘要正文……\par}
\end{minipage}

\vspace{6pt}
\hrule height 0.4pt
\vspace{10pt}

% ================= 正文 =================
\section{引言}

正文段落。首次出现图表时引用 \ref{fig:demo} 与表 \ref{tab:demo}。

\begin{figure}[htbp]
  \centering
  \includegraphics[width=0.86\textwidth]{figs/fig1.png}
  \caption{示例图题}
  \label{fig:demo}
\end{figure}

\begin{table}[htbp]
  \centering
  \caption{示例表题}
  \label{tab:demo}
  \footnotesize
  \begin{tabular}{@{}llr@{}}
    \toprule
    列一 & 列二 & 数值 \\
    \midrule
    A & 说明 & 1.00 \\
    \bottomrule
  \end{tabular}
\end{table}

\input{refs}

\end{document}
